\section{Multi-selector max-statistic validity}
\label{app:multiselector}

\paragraph{Claim.} Proposition~\ref{prop:multiselector-validity}.
\paragraph{Assumptions used.} A0.1--A0.5 (\cref{app:assumptions}).

Let $T^{(1)}_b, \dots, T^{(L)}_b$ denote the $L$ selector statistics computed on
permutation $b$ for a fixed feature, and define $M_b := \max_{\ell=1,\dots,L} T^{(\ell)}_b$.
If the joint vector of all selector statistics is exchangeable across
permutations under the null, then $(M_0, \dots, M_B)$ is exchangeable as a
measurable function of an exchangeable vector, so the right-tail +1 p-value computed from
$(M_b)$ is super-uniform by the same rank argument as in
\cref{them:plusone-superuniform}.

\begin{remark}[Scale/normalization]
The max-statistic construction is a validity statement. Practical power and
interpretability require care when mixing statistics on different scales. In
\citrees, selectors used together in multi-selector mode are chosen to be
comparable (e.g., bounded in $[0,1]$).
\end{remark}
